\section{总结与延伸}

\begin{enumerate}
    \item 从预训练 LM 到助手需要三步：instruction finetuning、奖励建模、策略优化。
    \item \textbf{RLHF}：训练奖励模型 + PPO 优化，是 ChatGPT 等模型的核心技术。
    \item \textbf{DPO}：将偏好优化简化为监督学习，无需 RL 基础设施。
    \item KL 正则化防止 reward hacking，是偏好优化的关键组件。
    \item RL 在 LLM 推理（数学、编程）中展现出巨大潜力。
    \item 偏好比较比绝对评分更可靠，Bradley-Terry 模型是标准框架。
    \item 部署阶段需要 evidence dashboard、drift alarm 与 ops playbook，形成实时监控闭环。
    \item 将每次 drift intervention 写入 runbook，与 governance checklist 联动，确保审计跟踪和快速恢复。
\end{enumerate}

\subsection{总结表}

\begin{table}[H]
\centering
\begin{tabular}{p{0.2\textwidth} p{0.35\textwidth} p{0.35\textwidth}}
\toprule
主题 & 核心收获 & 工程行动 \\
\midrule
Instruction Finetuning & 提供基本能力 & 用大规模 instruction data 初始化模型 \\
RLHF + KL & 通过奖励模型 + KL 限制实现 alignment & 维持 reference policy，避免 reward hacking \\
DPO & 偏好优化可以不依赖 RL & 直接训练策略，降低基础设施复杂度 \\
Pipeline Monitoring & 需要 timeline + checklist & 建立 dashboard、guardrail checklist \\
Evidence Matrix & 多信号联合验证对齐 & 把 benchmark, drift, red team 和 logs 结构化 \\
Deployment Trust & Pre-mortem/post-mortem + multi-role governance & 建立 rollback runbook、明确责任与 evidence matrix 对齐 \\
\bottomrule
\end{tabular}
\caption{RLHF 架构及执行要点}
\end{table}

\subsection{拓展阅读}
\begin{itemize}
    \item Ouyang et al., \textit{Training Language Models to Follow Instructions with Human Feedback} (InstructGPT, 2022)。
    \item Rafailov et al., \textit{Direct Preference Optimization: Your Language Model is Secretly a Reward Model} (DPO, 2023)。
    \item DeepSeek-AI, \textit{DeepSeek-R1: Incentivizing Reasoning Capability in LLMs via Reinforcement Learning} (2025)。
    \item Chung et al., \textit{Scaling Instruction-Finetuned Language Models} (FLAN-T5, 2022)。
\end{itemize}

\end{document}
